%==============================================================================
\section{Results}
\label{sec:results}

% --- Fill from the analysis pipeline (05_Execution_Plan.md, Tier 1 first).
%     Numbers come from generated/values.tex; [TODO] macros mark pending runs.

\subsection{Mutual information vs.\ $R^2$}
\label{sec:res-mi}
\Cref{tab:mi_vs_r2} sets the information measures beside the linear $R^2$
baseline on the PID campaign. The binned (Miller--Madow) estimate gives
$\MI{\slk}{\Tdie,\Vcore}=\valMIstv$~bits, with temperature dominating
($\MI{\slk}{\Tdie}=\valMIst$ vs.\ $\MI{\slk}{\Vcore}=\valMIsv$~bits); the KSG
estimator independently yields $\valKSGstv$~bits, and the moving-block bootstrap
95\,\% interval on the binned value is $[\valMIciLo,\valMIciHi]$~bits. The
cross-estimator agreement is the robustness evidence of \cref{sec:methodology}.

\paragraph{The non-linearity result.} The correct, apples-to-apples comparison is
\emph{not} the entropy fraction against the variance fraction (different scales),
but the measured MI against the \emph{Gaussian-equivalent} MI implied by the
linear fit, $-\tfrac12\log_2(1-R^2)=\valMIlinEquiv$~bits for
$R^2=\valRsqLinear$. The measured dependence ($\valKSGstv$~bits) \textbf{exceeds}
this linear-equivalent value by $\valMIexcessPct\,\%$: there is genuine
environmental information that a linear model cannot see---exactly the gap this
project set out to expose, and direct quantitative support for replacing $R^2$
with mutual information in this domain.

\paragraph{Surrogate null floor.} Because MI estimators are biased upward on
finite, autocorrelated samples, the measured value must be checked against the
floor the estimator produces under genuine independence. Circularly shifting the
$(\Tdie,\Vcore)$ stream against $\slk$ by large random lags---destroying the
cross-dependence while preserving each channel's marginal \emph{and}
autocorrelation---yields a null $\MI{\slk}{\Tdie,\Vcore}$ of mean
$\valMInullMean$~bits and 95th percentile $\valMInullFloor$~bits. The measured
$\valMIstv$~bits is $\approx\valMIfloorRatio\times$ this floor, so the dependence
is overwhelmingly real rather than estimator bias. The same floor frames the
post-correction reversible residual ($\valResidMI$~bits, \cref{sec:res-correction}):
having driven the leftover coupling down to the same order as the null floor is
exactly what ``the environment has been removed'' means in information terms.

\begin{figure}[t]
  \centering
  \includegraphics[width=0.86\linewidth]{fig_mi.pdf}
  \caption{Per-channel and joint mutual information between slack and the
  environment (bits). The joint $\MI{\slk}{\Tdie,\Vcore}$ exceeds the
  Gaussian-equivalent of the linear $R^2$ (dashed) by $26\%$, and sits well above
  the surrogate null floor (dotted), confirming genuine non-linear coupling.}
  \label{fig:mi}
\end{figure}

% Headline table: information measures vs. the linear R^2 baseline.
% [TODO] cells are filled by the Tier-1 MI analysis (see 05_Execution_Plan.md).
\begin{table}[t]
  \centering
  \caption{Information measures (bits) vs.\ the linear $R^2$ baseline on the PID
           campaign. The measured $\MI{\slk}{\Tdie,\Vcore}$ exceeds the
           Gaussian-equivalent MI of the linear fit
           ($-\tfrac12\log_2(1-R^2)$) by \valMIexcessPct\,\%, evidencing
           non-linear coupling. Prior-work columns are the published
           thermal-fidelity statistics.}
  \label{tab:mi_vs_r2}
  \begin{tabular}{lccc}
    \toprule
     & PID (new run) & PID (prior) & Bang-bang (prior) \\
    \midrule
    linear $R^2(\slk\mid\Tdie,\Vcore)$       & \valRsqLinear & \valRsqPid & \valRsqBangBang \\
    \;$\hookrightarrow$ Gaussian-equiv.\ MI (bits) & \valMIlinEquiv & --- & --- \\
    $\Ent{\slk}$ (bits)                      & \valHs        & ---        & --- \\
    $\MI{\slk}{\Tdie}$ (bits)                & \valMIst      & ---        & --- \\
    $\MI{\slk}{\Vcore}$ (bits)               & \valMIsv      & ---        & --- \\
    $\MI{\slk}{\Tdie,\Vcore}$ binned (bits)  & \valMIstv     & ---        & --- \\
    \;$\hookrightarrow$ KSG cross-check (bits) & \valKSGstv  & ---        & --- \\
    \;$\hookrightarrow$ 95\,\% CI (block btstrp) & [\valMIciLo,\,\valMIciHi] & --- & --- \\
    info.\ fraction $\MI{}{}/\Ent{}$         & \valInfoFraction & ---     & --- \\
    \bottomrule
  \end{tabular}
\end{table}


\subsection{PVT correction as information recovery}
\label{sec:res-correction}
We fit $\hat{\dpvt}(\Tdie,\Vcore)$ and form $\scorr=\slk-\hat{\dpvt}$. The
ordinary-least-squares (linear) correction drives the linear coupling to zero by
construction ($\operatorname{corr}(\scorr,\Tdie)=+0.00$,
$\operatorname{corr}(\scorr,\Vcore)=-0.00$); the methodologically relevant
question is how much \emph{reversible} mutual information it removes. Measured on
detrended signals---so that the slow aging drift does not masquerade as
environmental coupling---the reversible coupling falls from
$\valPreMIrev$~bits before correction to $\valResidMI$~bits after, i.e.\ the
linear correction removes \textbf{\valMIreductionPct\,\%} of the reversible
environment--sensor information, leaving a residual at the estimator's noise
floor. The reversible residual standard deviation is $\valResidSigmaLin$~counts,
consistent with the prior work's published PID residual ($0.73$~counts)---an
independent cross-check of the pipeline.

The physics-motivated non-linear correction (Arrhenius-like in $\Tdie$,
super-linear in $\Vcore$) removes marginally more residual MI
($\valResidMInl$ vs.\ $\valResidMIlin$~bits) but reduces the residual variance
negligibly; the MDL/BIC criterion (\cref{eq:mdl}, \cref{sec:th-mdl}) therefore
\textbf{selects the linear model} ($\valCorrModel$), judging the extra
parameters unjustified at this stress point---a description-length verdict
revisited in \cref{sec:discussion}.

\paragraph{Aging recovery (key result).} Removing the PVT contribution does not
merely clean the signal---it \emph{exposes} the latent aging trajectory. The
correlation of slack with time strengthens from
$\operatorname{corr}(\slk,\tm)=\valCorrStRaw$ on the raw stream to
$\operatorname{corr}(\scorr,\tm)=\valCorrStCorr$ after correction: stripping the
reversible environmental information makes the permanent aging drift the dominant
structure in the residual, which is precisely the information-recovery objective
of \cref{eq:model} and the empirical anchor for the RUL bound of
\cref{sec:res-rul}.

\begin{figure}[t]
  \centering
  \includegraphics[width=\linewidth]{fig_correction.pdf}
  \caption{(a) Raw slack tracks die temperature (corr $-0.86$). (b) After PVT
  correction the reversible coupling is gone and the permanent aging trend is
  \emph{exposed} (slack--time correlation strengthens from $-0.35$ to $-0.73$);
  the dashed line is the fitted aging rate.}
  \label{fig:correction}
\end{figure}

\subsection{Distributional distance between benches}
\label{sec:res-kl}
Two $\sim$24~h stabilized runs on different benches---the non-PID SBCCI setup
and the PID-controlled reference---let us put a number on ``how far the
non-compliant environment pushes the slack distribution.'' Both runs are on a
dedicated test device, distinct from the 214~h aging device of the preceding
sections; they share that test device, so the divergence isolates the
bench/environment difference rather than device-to-device variation. The benches
sit at different operating points, so the comparison is made on the centered
slack fluctuation, isolating noise character from setpoint. The slack noise is
$\valSlackStdSBCCI$ counts (SBCCI) vs.\ $\valSlackStdPID$ counts (PID), tracking
the temperature noise ($\valTstdSBCCI$ vs.\ $\valTstdPID$~$^\circ$C)---the
contamination is thermally driven.

The Kullback--Leibler divergence is strongly \emph{asymmetric}:
$\KL{p_{\text{SBCCI}}}{p_{\text{PID}}}=\valKLsp$~bits (the cost of describing the
noisy bench with the tight PID reference) versus
$\KL{p_{\text{PID}}}{p_{\text{SBCCI}}}=\valKLps$~bits---a direct illustration of
the course's point that KL is not a metric. The symmetric Jensen--Shannon
divergence is $\valJSbench$~bits (of a maximum of 1), a single transferable
``distance between validation environments.'' A zero-mean Gaussian closed form
gives $1.67$~bits for $\KL{p_{\text{SBCCI}}}{p_{\text{PID}}}$; the empirical
value is higher ($\valKLsp$~bits), and that excess is itself evidence that the
non-compliant bench injects \emph{non-Gaussian}, heavy-tailed excursions---which
motivates the blind-separation analysis of \cref{sec:res-ica}.

\begin{figure}[t]
  \centering
  \includegraphics[width=0.86\linewidth]{fig_kl.pdf}
  \caption{Centered slack-fluctuation distributions of the two benches (same test
  device). The non-PID SBCCI bench is markedly wider than the PID reference;
  the divergence $\KL{p_{\text{SBCCI}}}{p_{\text{PID}}}=\valKLsp$~bits
  (JS $=\valJSbench$~bits) quantifies the bench-fidelity gap.}
  \label{fig:kl}
\end{figure}

\subsection{Effective resolution and alarm reliability}
\label{sec:res-capacity}
Taking the post-correction reversible noise $\sigma_{\text{noise}}=\valSigNoise$
counts and the aging-signal spread $\sigma_{\text{sig}}=\valSigSignal$ counts,
the capacity heuristic (\cref{eq:capacity}, derived in \cref{sec:th-capacity}
from the Gaussian maximum-entropy argument) gives $\mathrm{SNR}=\valSNR$,
$C=\valCapacityBits$~bits and therefore $\valCapacityStates$ distinguishable
aging states over this campaign---a deliberately modest figure that reflects an
aging excursion only $\sim\!\sqrt{\valSNR}\times$ the noise floor, and that the
prior residual-$\sigma$ metric could not express.

The aging-alarm decision is modelled as a BSC whose crossover $p$ follows from
the per-window detection SNR (\cref{eq:bsc}). Reliability is strongly
\emph{time-scale dependent}: a per-hour decision is essentially a coin flip
($p\approx0.49$, $\Cbsc\approx0$) because one hour of drift is buried in the
noise, whereas a $\valBSCwindow$~h operational window gives $p=\valBSCp$ and
$\Cbsc=\valBSCcap$~bits/decision (near-perfect), and the full campaign is
effectively certain. The crossover between these regimes coincides with the
minimum observation time derived next.

\begin{figure}[t]
  \centering
  \includegraphics[width=0.82\linewidth]{fig_capacity_rul.pdf}
  \caption{Aging-alarm channel capacity $\Cbsc=1-\Hb(p)$ as a function of the
  decision-window length. Reliability climbs from $\approx0$ (per-hour, buried in
  noise) to $\approx1$ bit, with the transition coinciding with the
  Cram\'er--Rao minimum observation time ($\approx\valMinObsTime$~h, dashed)---the
  three figures of merit are mutually consistent.}
  \label{fig:capacity}
\end{figure}

\subsection{RUL detectability bound}
\label{sec:res-rul}
Posing RUL as least-squares estimation of the aging slope, and using the measured
noise with the effective sample size $\neff=\valNeff$
($\tauint=\valTauInt$~samples, i.e.\ $\sim$5~min decorrelation), the
Cram\'er--Rao bound (\cref{sec:th-estimation}) gives a slope-estimate standard deviation of $\valCRLBsd$
m-cnt/h over the $\sim$214~h campaign, hence a minimum detectable aging rate
(at $3\sigma$) of \textbf{\valMinDetectRate~m-cnt/h}. The measured corrected
aging rate is $\valAgingSlope$~m-cnt/h---roughly $15\times$ above this floor, so
the trend is firmly detectable. Equivalently, the minimum observation time to
resolve the measured trend at $3\sigma$ is \textbf{$\sim$\valMinObsTime~h}, well
within this campaign (and explaining why the prior $\sim$23~h runs could not
separate aging from the reversible floor). This is the honest detectability
statement the project promised: not a lifetime prediction, but a noise-limited
bound on \emph{when} aging becomes observable.

\paragraph{Bayesian credible bars.} The frequentist bound is complemented by a
Bayesian linear regression of $\scorr$ on time, fit on the data thinned to the
$\neff$ effectively-independent points (so the i.i.d.\ likelihood holds) under a
Jeffreys prior, giving a Student-$t$ posterior for the slope. The posterior aging
rate is $\valBayesSlope$~m-cnt/h with a $95\%$ credible interval of
$[\valBayesSlopeHi,\valBayesSlopeLo]$~m-cnt/h, and the posterior probability that
the slope is genuinely negative---that aging is present---is
$\valBayesPaging\,\%$, numerically indistinguishable from certainty. Propagating
the slope posterior through the detectability relation gives a credible interval
of $[\valBayesMotLo,\valBayesMotHi]$~h on the minimum observation time, bracketing
the $\valMinObsTime$~h point estimate. The two paradigms agree, and the Bayesian
treatment supplies the calibrated error bars a future RUL study would propagate
forward.

\subsection{A denoised aging trajectory and RUL as a distribution (Wiener/Kalman)}
\label{sec:res-wiener}
The Cram\'er--Rao and Bayesian results above give a single campaign-wide rate
and its error bar. Here we ask the sequential-estimation question of
\cref{sec:th-wiener} instead: block-averaging $\scorr$ into $\valNBlocks$
quasi-independent observations (one per $2\tauint$, i.e.\ one $\neff$ unit
each), the method-of-moments fit of \cref{eq:wiener-model} gives a Wiener
diffusion $\sigma_B=\valWienerSigmaB$ counts/$\sqrt{\text{h}}$ and a
block-mean observation noise $R$ notably \emph{below} the native
per-sample $\sigma_{\text{noise}}^2$ (\cref{sec:res-capacity}) --- direct
evidence that averaging over a full decorrelation window still removes
sub-$\tauint$ structure that the single autocorrelation time alone does not
capture. The rate-noise MLE collapses to its numerical floor
($\sigma_{\mathrm{rate}}^2\!\approx\!\valWienerQRate$): the data prefer a
\emph{constant}-drift Wiener process over a time-varying one, a second,
independent confirmation --- alongside the MDL verdict of
\cref{sec:res-correction} --- that this single, narrow-operating-point
campaign does not contain enough curvature information to support a more
flexible model. The Kalman/RTS-smoothed rate at the end of the campaign is
$\valWienerRateEnd$~m-cnt/h, consistent in sign and order of magnitude with
the CRLB/Bayesian estimate of \cref{sec:res-rul} ($\sim\!-18$~m-cnt/h); the
$\sim$15\% gap between the two independently-derived rates reflects the
different estimators (whole-series OLS vs.\ a filtered, block-averaged
state), not a contradiction. The smoothed level's posterior standard
deviation is essentially flat across the campaign
($\valWienerLevelSDstart\to\valWienerLevelSDend$ counts): the filter reaches
its steady-state uncertainty within the first few blocks and does not
keep shrinking indefinitely, exactly as a linear-Gaussian filter with
constant $R$ and $\sigma_B^2$ predicts.

\begin{figure}[t]
  \centering
  \includegraphics[width=\linewidth]{fig_wiener_rul.pdf}
  \caption{(a) Block-mean $\scorr$ (light) and the Kalman/RTS-smoothed aging
  trajectory $\hat\Xlat(t)$ with a 95\% credible band. (b) The Inverse-Gaussian
  RUL distribution (\cref{eq:invgauss}) at two Tier-1-anchored thresholds: one
  noise-floor increment ($D=\sigma_{\text{noise}}$) and one
  already-observed signal-scale excursion ($D=\sigma_{\text{signal}}$).}
  \label{fig:wiener}
\end{figure}

\paragraph{RUL as a distribution.} At $D=\sigma_{\text{noise}}=\valDnoise$
counts (the same noise floor that limits the sensor to
\valCapacityStates{} distinguishable states, \cref{sec:res-capacity}), the
Wiener first-passage time has mean $\valRULmeanNoise$~h but median
$\valRULmedianNoise$~h; at $D=\sigma_{\text{signal}}=\valDsignal$ counts, mean
$\valRULmeanSignal$~h and median $\valRULmedianSignal$~h. The large gap
between mean and median in both cases is not an artifact: the Inverse
Gaussian is strongly right-skewed exactly when the threshold $D$ is small
relative to the diffusion scale $\sigma_B$, which is precisely the
low-SNR ($\mathrm{SNR}=\valSNR$, \cref{sec:res-capacity}) regime measured
here --- \emph{the same noise limitation that caps the sensor's resolvable
states also makes short-horizon RUL predictions highly uncertain}, a
quantitative link between the capacity result and the RUL result that a
single point estimate cannot express. This is offered as a
\emph{methodology} for a full RUL distribution, exercised at self-referential
thresholds because, as throughout this paper, no externally calibrated
failure threshold exists for this sensor; supplying one (\cref{sec:conclusion})
turns \cref{eq:invgauss} directly into an operational RUL distribution.

\subsection{Blind source separation (ICA)}
\label{sec:res-ica}
As a model-free counterpart to the regression correction of
\cref{sec:res-correction}, FastICA (negentropy contrast, \cref{eq:negentropy} in
\cref{sec:th-ica}) was applied to the three
observed channels $(\slk,\Tdie,\Vcore)$ alone, with no functional form assumed
for $\dpvt$. The recovered component that tracks time matches the
regression-based aging residual $\scorr$ at
$|\mathrm{corr}|=\valICAagingScorr$---\emph{model-free ICA independently arrives
at essentially the same aging signal the regression produced}, which
cross-validates the PVT correction without using the measured $\Tdie,\Vcore$
coupling. A second component aligns with temperature
($|\mathrm{corr}(\cdot,\Tdie)|=\valICAthermalT$). Both recovered sources are
strongly non-Gaussian (excess kurtosis $\valICAkurtAging$ and
$\valICAkurtThermal$), confirming the premise that non-Gaussianity is the feature
distinguishing the physical sources---and consistent with the heavy-tailed
excursions the KL analysis (\cref{sec:res-kl}) exposed.

The separation is \emph{partial}, and is reported as such: the residual pairwise
mutual information between recovered components reaches $\valICAresidMI$~bits
(ideal independence is $0$), because the slow, non-stationary aging drift and the
heavy-tailed noise only approximately satisfy ICA's i.i.d./stationarity
assumptions. ICA is therefore corroborative here---an independent, assumption-light
route to the same aging component---rather than a clean unmixing. That a method
using none of the $\Tdie,\Vcore$ coupling still recovers $83\%$ of the regression
aging signal is itself the result.

\begin{figure}[t]
  \centering
  \includegraphics[width=\linewidth]{fig_ica.pdf}
  \caption{The ICA-recovered aging component (purple) overlaid on the
  regression-based residual $\scorr$ (green), both $z$-scored. Model-free ICA,
  given only $(\slk,\Tdie,\Vcore)$, tracks the regression aging signal at
  $|\mathrm{corr}|=\valICAagingScorr$.}
  \label{fig:ica}
\end{figure}
